

%% LyX 2.4.4 created this file.  For more info, see https://www.lyx.org/.
%% Do not edit unless you really know what you are doing.
\documentclass[12pt,onecolumn]{paper}
\usepackage{amsmath}
\usepackage{amsthm}
\usepackage{newtxmath}
%\usepackage[latin9]{inputenc}
\usepackage[utf8]{inputenc}
\usepackage[T1]{fontenc}
\usepackage{color}
\usepackage{array}
\usepackage{verbatim}
\usepackage{mathrsfs}
\usepackage{multirow}
\usepackage{graphicx}
\usepackage{geometry}
\geometry{verbose}
\usepackage[bookmarks=false,
 breaklinks=false,pdfborder={0 0 1},backref=false,colorlinks=true]{hyperref}
 %\usepackage[utf8]{inputenc}  % (not needed in modern LaTeX, but sometimes helps)
%\usepackage[T1]{fontenc}
\makeatletter

%%%%%%%%%%%%%%%%%%%%%%%%%%%%%% LyX specific LaTeX commands.
%% Because html converters don't know tabularnewline
\providecommand{\tabularnewline}{\\}

%%%%%%%%%%%%%%%%%%%%%%%%%%%%%% Textclass specific LaTeX commands.
% Save and disable \example as this might
% clash with theorems
\let\paperclassexample\example
\let\endpaperclassexample\endexample
\let\example\relax
% If theorems hasn't been loaded, restore example
\AtBeginDocument{%
  \@ifundefined{example}{\let\example\paperclassexample}{}
}

%%%%%%%%%%%%%%%%%%%%%%%%%%%%%% User specified LaTeX commands.
\usepackage{amsfonts}
\usepackage{newtxtext}
\usepackage{enumitem}
\usepackage{titling}


\setlength{\droptitle}{-6em}

% Enter the specific assignment number and topic of that assignment below, and replace "Your Name" with your actual name.
%\title{SAXS super resolution status}
%\author{Uri Steinitz}
%\date{Sep. 2022}


% keep these (hide display, keep data for linking)
%\usepackage[bibstyle=numeric-comp,citestyle=numeric,doi=false,url=false,eprint=false]{biblatex}
\usepackage[style=nature,backend=biber,
  giveninits=true,terseinits=true,
  maxbibnames=6,minbibnames=1,
  doi=false,url=false,eprint=false]{biblatex}

\addbibresource{saxs2.bib}

% hide the 'note' field
\AtEveryBibitem{\clearfield{note}}
\AtEveryCitekey{\clearfield{note}}

% Clickable (year): URL > DOI > arXiv(eprinttype=arxiv); always parentheses
\renewbibmacro*{date}{%
  \iffieldundef{year}
    {}
    {\printtext{%
        \iffieldundef{url}
          {\iffieldundef{doi}
            {\iffieldundef{eprint}
              {\printfield{year}}
              {\iffieldequalstr{eprinttype}{arxiv}
                 {\href{https://arxiv.org/abs/\thefield{eprint}}{\printfield{year}}}
                 {\printfield{year}}}}
            {\href{https://doi.org/\thefield{doi}}{\printfield{year}}}}
          {\href{\thefield{url}}{\printfield{year}}}%
      }%
    }%
}
%\renewbibmacro*{date+extrayear}{%
%  \usebibmacro{date}% (already prints parentheses)
%  \setunit*{\addspace}\printfield{extrayear}% e.g., 2007a/2007b
%}

% Prefer shortjournal if present; otherwise use full journaltitle
\DeclareFieldFormat{journaltitle}{%
  \mkbibemph{\iffieldundef{shortjournal}{#1}{\printfield{shortjournal}}}%
}

\begin{document}
\include{saxs2.bib}
\title{Variable-Resolution Scattering Reveals Ensemble Properties}
\author{Uri Steinitz, Roy Beck}

\maketitle
\begin{abstract}
Small-angle X-ray and neutron scattering (SAXS/SANS) provides a commonplace characterization technique for a wide variety of nanoscopic systems. In practice, the nanoscopic scale presents an inherent challenge as many studied samples exhibit a distribution of structural properties (polydispersity) where this distribution influences their function. In addition, structural polydispersity convolute the SAXS/SANS intensity profiles in a similar manner to that of the instrumental induced point-spread function.    
Here, we introduce a practical method (\textbf{TENOR-SAXS}) that leverages changes in resolution to uncover the ensemble's variance directly from 2D scattering data, without any \textit{a priori} knowledge on the sample geometry or instrument resolution.   By employing a data-efficient comparison approach, we demonstrate that ensemble variance can be reliably estimated for synchrotron-scale noise-levels. We show that the information can be retrieved through digital smearing of a single measured image, eliminating the need for additional experimental measurements which can be applied to already previously collected SAXS/SANS data.

\end{abstract}

\section{Introduction}

Small-angle X-ray scattering (SAXS) is a widely used technique for characterizing nanostructures, providing ensemble-averaged information on particle size, shape, and conformational flexibility. Small angle neutron scattering (SANS) shares similar analysis techniques, though its availability, resolution and noise level are not as superior as SAXS'. Since the early developments of Guinier and co-workers \cite{guinier1955,glatter1982}, SAXS has become central in nanoscopic characterization, in particular for samples measured in solutions. A relevant example includes the study of nanoscale biological substances (i.\,e., proteins, lipids, nucleic acids and complexes of those), where it complements high resolution crystallography and NMR by accessing macromolecular length-scale and states in solution, including flexible and heterogeneous systems \cite{svergun2003,jacques2010}. In soft condensed matter and materials science, SAXS is similarly applied to polymers, nano-particles, colloids, and porous media, enabling quantitative description of morphology and hierarchical structure \cite{feigin1987,pauw2013}.

The distribution of nano-particle sizes has fundamental effects in many cases in materials sciences and pharmaceuticals \cite{schoepe2007,garcia2009,pastoriza2011,danaei2018}. This prompts for devising SAXS methods to observe and characterize the polydispersity of a nanoscopic ensemble.  SAXS is fundamentally limited, since the measured intensity represents an average over all scatterers in the illuminated volume. Samples such as intrinsically disordered proteins (IDP) exhibit structural dispersion, which may be effectively represented in the measured profile as a distribution of the radius of gyration ($R_g$), rather than a single structure \cite{ahmed2020}. The average value that conventional SAXS analysis often measure not only masks this dispersion but is even biased by it. Approaches to recover ensemble heterogeneity include the ensemble optimization method \cite{bernado2007,tria2015}, Bayesian inference \cite{antonov2016,chatzimagas2023}, molecular form factor \cite{riback2017} and related modeling frameworks. Recently, Tomchuk et al. introduced a new analysis protocol to extract polydispersity directly from SAXS \cite{tomchuk2023} by extracting observable values from the data. However, this analysis is sensitive to the exact functional form of the $R_g$ distribution, which is unknown in most cases, and assumes spherical particles.      

Instrumental effects provide an additional complication. Beam divergence and detector resolution broaden the measured scattering profile through convolution with the point-spread function (PSF, or resolution function). This broadening mimics the effect of particle-size polydispersity, so that accurate recovery of ensemble variance from a single scattering profile requires knowledge of the PSF. However, PSF characterization is often incomplete and may vary with beamline configuration \cite{pedersen1990,vad2010,rennie2013}. Although deconvolution approaches are sometimes used in principle to correct for instrumental broadening, deconvolution itself requires accurate knowledge of the PSF and is typically ill-posed; as a result, it propagates PSF uncertainty directly into the inferred variance and can introduce systematic bias \cite{tikhonov1977,bertero1998}.

In this work, we introduce a Technique for ENsemble Observation by Resolution variation in SAXS (\textbf{TENOR-SAXS}), which eliminates the dependence on PSF characterization and avoids the need for deconvolution. The method leverages changes in instrumental resolution to construct observables whose dependence on the variance of $R_g$ is analytically defined and is rather robust against the specific instrumental PSF and the exact probability distribution function of $R_g$. Scattering patterns acquired at different effective resolutions, either experimentally (e.g.\ via slit or collimation adjustments) or by applying digital smearing, are compared to separate instrumental and structural contributions to profile broadening. In this way, the ensemble's polydispersity is characterized without requiring explicit knowledge of the PSF and without the systematic bias associated with deconvolution-based approaches.

Below, we present the theoretical formulation of the approach, validate it using analytical models, and demonstrate through simulation that ensemble variance can be recovered with minimal assumptions concerning the underlying monodisperse form factor. We further evaluate the effects of experimental noise and show that signal levels typical of modern synchrotron SAXS measurements are compatible with reliable implementation.

These results establish TENOR-SAXS as a practical means of quantifying structural dispersion in solution-based ensembles, while eliminating PSF-induced bias in estimates of variance. Currently, the method works best when the monodisperse form-factor is known, and it has only a minimal dependence on the particle size distribution function. As such, it is applicable to ensembles of nano-scale assemblies with predetermined shapes. Future work will aim to use the approach on particles with less distinct shapes such as disordered proteins and polymers.

\section{Theory}
\subsection{Scattering off a polydispersive ensemble}
\label{sec:scattering-polydisp}
In scattering experiments (i.\,e., SAXS/ SANS), the scattered intensity is measured as a function of the scattering vector $q = \frac{4\pi}{\lambda}\sin\theta$, where $\lambda$ is the wavelength of the X-rays or neutrons and $2\theta$ is the scattering angle. In an ideal setting of zero instrumental smearing, the measured intensity $I(q)$ reflects an average over all particles in the illuminated volume. For a system of identical particles (monodisperse), this intensity is directly related to the single particle form factor. However, as recognized by Guinier, many samples are polydisperse, meaning that structural properties such as the $R_g$ are distributed around a mean value rather than being fixed \cite{guinier1955}.   

Therefore, for heterogeneous systems, the observed (normalized) scattering can be written as an ensemble average
\begin{equation}
I(q) = \big\langle F(q, R_g) \big\rangle ,
\label{eq:ensemble_average}
\end{equation}
where $F(q, R_g)$ is the scattering intensity from a particle of size $R_g$, and the brackets denote averaging over the distribution of $R_g$ values in the ensemble. We note here that the form factor averaging is weighted, including both the numerical abundance of each $R_g$ and the scattering strength. In the analysis presented here we use numerical abundance only, so the scattering strength is implemented by altering the distribution function. To quantify the spread of this distribution, we define its scattering weighted mean as $R_0$ and introduce a dimensionless variance
\begin{equation}
V = \frac{\mathrm{Var}(R_g)}{R_0^2},
\end{equation}
which uses, as with the mean value, the scattering strength weighted distribution. If the distribution is narrow ($V\ll 1)$, we can expand $\ln F(q,R_g)$ around $R_0$, to leading order in $V$:

\begin{equation}
\ln I(q) \;\approx\;\ln I_0\,+\, \ln F(q, R_0) \,+\, \tfrac{1}{2} V R_0^2 \,\frac{\partial^2}{\partial R_g^2} \ln F(q, R_g)\Big|_{R_g = R_0}.
\label{eq:variance_correction}
\end{equation}

Examples of ensembles featuring $R_g$ distributions include estimates of $V$ spanning values of 0.01-0.25 for different samples \cite{ullman2011,Kikhney2015,ahmed2020,pauw2017}. At low values of $qR_0$ the Guinier approximation is valid, and the scattering of a monodisperse particle of size $R_g$ takes the form (to first order of $q^2 R_g^2$):
\begin{equation}
F(q,R_g) \approx I_0\exp\!\left(-\tfrac{1}{3} q^2 R_g^2\right).
\label{taylor-F}
\end{equation}
Substituting Eq.\,\eqref{taylor-F} into Eq.~\eqref{eq:ensemble_average} and retaining terms to leading order in $V$ yields
\begin{equation}
\ln I(q) \;\approx\; \ln I_0 - \tfrac{1}{3} q^2 R_0^2 \,(1 + V) = \ln I_0 - \tfrac{1}{3} q^2 \tilde{R}_0^2 .
\label{eq:guinier_variance}
\end{equation}

Equation~\eqref{eq:guinier_variance} shows that polydispersity manifests as an effective increase in the apparent radius of gyration when analyzed in the Guinier regime. While a monodisperse system would yield a slope (of the curve $\ln I (q^2)$) proportional to $R_0^2$, the additional factor $(1+V)$ encodes the distribution of particle sizes (for any distribution function). Guinier analysis \emph{per se} does not distinguish between a monodisperse sample with a larger radius of gyration and an ensemble with a distribution of radii. Thus, we develop a more elaborate method to estimate the variance, from which also the mean radius of gyration could be extracted. The concept uses the effect that the convolution of the intrinsic scattering with different instrumental point-spread functions yields different intensity maps. The comparison of such data, obtained for example by measurements with varying slit sizes, encodes also details of the ensemble's polydispersity.


\subsection{Resolution effects and convolution with the instrumental point-spread function}


\label{sec:psf-smearing}

Let us consider a monodisperse sample where the measured intensity is the convolution of the intrinsic form factor \(F(\mathbf q)\) with the instrument point-spread function (PSF):  \cite{guinier1955,pedersen1990,pedersen1997}:
\begin{equation}
I(\mathbf q) \;=\; R (q) \circledast F(q) \;=\; \int R(-\mathbf p)\,F(\mathbf q+\mathbf p)\,d^2p ,
\label{eq:conv-main}
\end{equation}
where $R(q)$ represents the instrumental PSF.   
In typical small-angle setups, rectangular aperture slits and anistropic beam profiles make \(R(q)\) well approximated by an anisotropic Gaussian \cite{vad2010},
\begin{equation}
R(\mathbf q) \;=\; \frac{1}{2\pi\sigma_x\sigma_y}\,
\exp\!\Big(-\frac{q_x^2}{2\sigma_x^2}-\frac{q_y^2}{2\sigma_y^2}\Big),
\label{eq:R-gauss}
\end{equation}
with widths \(\sigma_x\) and \(\sigma_y\) along the aperture axes, assumed to be aligned with the detector.

When the PSF widths are small compared to the second $q$ log-derivative (typically $ \sigma^2\cdot R_0^2\ll1 $), a second-order expansion of the convolution gives the following expression for the log-intensity (see Appendix \ref{app:direct-logI} for derivation):
\small{
\begin{equation}
\ln I(\mathbf q)\;\approx\; \ln F(q)
\;+\;\tfrac{1}{2}\big(\sigma_x^2\,\partial_{q_x}^2+\sigma_y^2\,\partial_{q_y}^2\big)\ln F( q)
\;+\;\tfrac{1}{2}\Big[\sigma_x^2\big(\partial_{q_x}\ln F( q)\big)^2+\sigma_y^2\big(\partial_{q_y}\ln F( q)\big)^2\Big]\
\label{eq:logI-operator}
\end{equation}
}
For samples with isotropic form factors ($F(\mathbf q)=F(|\mathbf q|))$, either by having spherical symmetry or due to angular ensemble averaging, Eq.~\eqref{eq:logI-operator} shows that the PSF-smearing correction to the log-intensity depends on the radial derivative of $\ln F$, with anisotropic weights set by \(\sigma_x^2\) and \(\sigma_y^2\).
Equation \eqref{eq:logI-operator} also demonstrates the resemblance between the effect of PSF and the polydispersity (Eq.\,\eqref{eq:variance_correction}). This work shows how to decouple the two effects with only a limited \emph{a priori} knowledge.


\subsection{Extraction of the variance from 2D measurements}
\label{sec:extraction2D}

We consider two measurements of the same sample taken with different resolutions characterized by anisotropic Gaussian widths $(\sigma_{x,1},\sigma_{y,1})$ and $(\sigma_{x,2},\sigma_{y,2})$. These PSFs may be realized either by taking two measurements with different settings (e.\,g., distances or slit sizes) or by smearing a single measurement using two different filters, eliminating the need for normalization. 
Let $F(q)$ be the isotropic intrinsic form-factor intensity (single-particle, no interparticle correlations) and $I_k(\mathbf q)$ the measured intensity for experiment $k\in\{1,2\}$, and the PSF width differences as:
\begin{equation}
\Delta\sigma_x^2=\sigma_{x,1}^2-\sigma_{x,2}^2,\qquad
\Delta\sigma_y^2=\sigma_{y,1}^2-\sigma_{y,2}^2.
\end{equation}
Using Eq.\,\eqref{eq:logI-operator}, the log-ratio of the two intensities is:
\begin{equation}
\begin{aligned}
\mathcal{G}(\mathbf q)&\;\equiv\;\ln I_1(\mathbf q)-\ln I_2(\mathbf q)
\\
&\;=\;\frac{1}{2}\,\Delta\sigma_x^2\Big[\partial_{q_x}^2\ln F+(\partial_{q_x}\ln F)^2\Big]
+ \frac{1}{2}\,\Delta\sigma_y^2\Big[\partial_{q_y}^2\ln F+(\partial_{q_y}\ln F)^2\Big].
\label{eq:G_cart}
\end{aligned}
\end{equation}

Given $q=\sqrt{q_x^2+q_y^2}$ on the detector plane, and the in-plane azimuthal angle $\chi$ (\textit{i.\,e.}, $q_x=q\cos\chi, q_y=q\sin\chi$), we can substitute the Cartesian derivatives in  Eq.~\eqref{eq:G_cart} to:
\begin{equation}
\mathcal{G}(q,\chi)\;=\;\frac{1}{2}\,\Sigma_0\Big[f'(q)+f(q)^2+\frac{f(q)}{q}\Big]
+\frac{1}{2}\,\Delta_0\Big[f'(q)+f(q)^2-\frac{f(q)}{q}\Big]\cos 2\chi,
\label{eq:G_polar_compact}
\end{equation}
where,  
\begin{equation}
\begin{aligned}
f(q)\equiv\frac{d}{dq}\ln F(q),\qquad \Sigma_0\equiv\tfrac{1}{2}(\Delta\sigma_x^2+\Delta\sigma_y^2),  \\
f'(q)\equiv\frac{d}{dq}f(q),\qquad
\Delta_0\equiv\tfrac{1}{2}(\Delta\sigma_x^2-\Delta\sigma_y^2).
\end{aligned}
\end{equation}

For compactness we define the radial functions:
\begin{equation}
G(q)\;=\;f'(q)+f(q)^2+\frac{f(q)}{q},\qquad
M(q)\;=\;f'(q)+f(q)^2-\frac{f(q)}{q},
\label{eq:GM_defs}
\end{equation}
and then Eq.~\eqref{eq:G_polar_compact} becomes:
\begin{equation}
\,\mathcal{G}(q,\chi)\;=\;\tfrac{1}{2}\,\Sigma_0\,G(q)\;+\;\tfrac{1}{2}\,\Delta_0\,M(q)\cos 2\chi\, .
\label{eq:G_final}
\end{equation}

Equation~\eqref{eq:G_final}, that is fully derived in Appendix \ref{app:direct-logI}, shows that the dependence on the two experimental resolution settings is only experienced through the pre-factors $\Sigma_0$ and $\Delta_0$ representing the PSFs' isotropic and anisotropic differences, respectively. 
The azimuth-independent contribution is proportional to $G(q)$, while the twofold anisotropic contribution is proportional to $M(q)\cos 2\chi$. 
Changing the PSF widths between experiments simply rescales these contributions, while the intrinsic functions $G(q)$ and $M(q)$ contain the raw ensemble information.


If we take a monodisperse particle with a specific $R_g$, it is natural to express its form-factor through a unit-less parameters $u=(qR_g)^2$. Following, we can expand the form-factor logarithm around the origin:
\begin{equation}
\ln F_{R_g}(q) = \phi(u) = \kappa\,u + \tfrac{1}{2}\phi''\,u^2 + O(u^3)    
\end{equation}
with
\begin{equation}
\phi(u)= \kappa\,u + \tfrac{1}{2}\phi''\,u^2 + O(u^3),
\qquad
\kappa=\partial_u\phi(0),\;\; \phi''=\partial^2\phi(0)/\partial u^2\;.
\label{eq:phi_series u}
\end{equation}
At low scattering angles, the Guinier rule applies\cite{guinier1955}, and $\kappa=-\tfrac{1}{3}$, and we limit our current discussion to the second order in $u$. In Tab.\,\ref{tab:Typical-values-for} we show some examples for typical form-factors but one may readily derive the relevant $\phi''$ from any specific form-factor, by taking the Maclaurin series of $\ln{F}$ with respect to the $u$. These examples demonstrate and partly justify the assumption that $\phi''\ll1$ for practical cases.

\begin{table}
    \centering
    \begin{tabular}{|c|c|c|c|}\hline
         Monodisperse form factor&  $\kappa$&  $\phi''$& Source\\\hline\hline
         Spherical shell& $-1/3$ & -1/45 &{\cite{guinier1955}{}} \\\hline
         Solid sphere&  $-1/3$&  -1/63& {\cite{rayleigh1910}}\\\hline
         Thin rod&  $-1/3$& +11/225& {\cite{pedersen1997}}\\\hline
         Gaussian chain&  $-1/3$& +1/18& {\cite{debye1947}}\\ \hline
    \end{tabular}
    \caption{Canonical monodisperse form-factors and their relevant $\kappa,\phi''$ coefficients (defined
by the unit-less derivatives in Eq.\,\eqref{eq:phi_series u}) \label{tab:Typical-values-for}}

\end{table}

For a collection of particles sharing an identical form-factor functional dependence (\textit{i.e.}\,a fixed $\phi''$), but exhibiting a distribution of radii of gyration $P(R_g)$, the polydisperse scattering intensity $I(q)$ can be expressed using a cumulant expansion. This method was originally introduced by Thiele \cite{thiele1889} and developed in modern form by Cram\'er \cite{cramer1946}, and has previously been applied to size polydispersity in scattering problems \cite{koppel1972} (see Appendix~\ref{app:cumulant}).

The expansion reaches a polynomial form of the functions $G$ and $M$: 
\begin{equation}
G(q)=\sum_{n\ge0} g_n (q^2)^n,\qquad
M(q)=\sum_{n\ge0} m_n (q^2)^n.
\label{eq:GM-poly}
\end{equation}
The coefficients $g_n$ and $m_n$ depend on $R_0$, $V$, and $\phi''$ (the radius of gyration's weighted mean and relative variance and the monodisperse form-factor curvature, respectively), with the approximate values presented in Tab.\,\ref{tab:GM-coef}. 
\begin{table}
     
\begin{center}
\begin{tabular}{|c|l|l|}
\hline
$\mathbf{n}$ & \textbf{$G$ coefficients ($g_n$)}& \textbf{$M$ coefficients ($m_n$)}\\
\hline\hline
0& $-\dfrac{4}{3}\,R_0^2(1+V) = -\dfrac{4}{3}\,\tilde{R}_0^2$ 
& $0$ 
\\[6pt]\hline
1& $R_0^4\!\cdot\, \dfrac{4}{9}\left(1+18\phi'' + \!10V+ 108V\phi''\right)$& $R_0^4\!\cdot\, \dfrac{4}{9}\left(1+9\phi'' + \!6V+ 54V\phi''\right)$\\[6pt]\hline
2& $-R_0^6\!\cdot\,\dfrac{8}{3}\left(\phi'' + \!\dfrac{4}{9}V+ 50V\phi''\right)$& $-R_0^6\!\cdot\, \dfrac{8}{3}\left(\phi'' + \!\dfrac{4}{9}V+ 19V\phi''\right)$\\[6pt] \hline
\end{tabular}
    \caption{Analytical expressions for the coefficients for the $G$ and $M$ functions, up to first order in $V$ and $\phi''$.}
    \label{tab:GM-coef}

\end{center}
\end{table}

\subsection{2D extraction of the functional form}
Equation~\eqref{eq:G_final} allows an effective separation of variables once we have a 2D log-ratio of two measurements. It is possible to fit (e.\,g., by the least-squares method) the functional form of the map of $\mathcal{G}$ to an azimuth-independent part and a twofold anisotropic component. Using this fitting scheme is data-efficient since it avoids binning and gains information directly from the full set of the raw data.

Fitting the map to the polynomial forms of $G$ and $M$ (Eq.\,\eqref{eq:GM-poly} and Tab.\,\ref{tab:GM-coef}) would leave a scaling parameter (seen in Eq.\,\eqref{eq:G_final}) which depends on the different PSFs used. However, it is possible to extract scale-independent information by using coefficient \emph{ratios}, such as $g_1/g_0$ or $m_2/m_1$. 

Using Eq.\,\eqref{eq:guinier_variance} and a Guinier analysis we can estimate $\tilde{R}_0^2=R_0^2\,(1+V)$. Therefore we may define unit-less observables which all originated from the data, such as:
\begin{equation}
\label{eq:observables-old}
\mathcal{J}_G \equiv \frac{g_1/g_0}{\tilde{R}_0^2}\approx\frac{\,1+18\phi'' + 10V\!+108V\phi''}{-3\,(1+V)^2}
\qquad,
\end{equation}

\begin{equation}
\mathcal{J}_{G2} \equiv \frac{g_2/g_1}{\tilde{R}_0^2}\approx\frac{18\phi''+8V+450\,V\phi''}{\,-3(1+V)(1+18\phi'' + 10V+108\,V\phi'')}
\qquad,
\end{equation}

\begin{equation}
\mathcal{J}_M \equiv \frac{m_2/m_1}{\tilde{R}_0^2}
\approx \frac{18\phi'' + 8 V +\! 342\,V\phi''}
{-3(1+V)\left(1+9\phi'' + 6V\!+54\,V\phi''\right)}.\label{eq:R_m-observe-old}
\end{equation}
Equations \eqref{eq:observables-old}-\eqref{eq:R_m-observe-old} present three unit-less independent observables from which $R_0$, the PSFs, and the beams' strengths were eliminated. Below, we will show how to extract the two remaining parameters, $V$ and $\phi''$ from these observables. We note that instead of dividing by the observed radius of gyration squared ($\tilde{R}_0^2$) , division by $g_0$ provides expressions which are $R_0$ independent:

\begin{equation}
\label{eq:observables}\mathcal{Y}_{G} \equiv \frac{g_1}{g_0^2}\approx\frac{\,1+18\phi'' +10V\!+108V\phi''}{4\,(1+V)^2}
\qquad,
\end{equation}

\begin{equation}
\mathcal{Y}_{G2} \equiv \frac{g_2}{g_1g_0}\approx\frac{18\phi''+8V+450V\phi''}{\,4(1+V)(1+18\phi'' + 10V+108V\phi'')}
\qquad,
\end{equation}

\begin{equation}
\mathcal{Y}_M \equiv \frac{m_2}{m_1g_0}
\approx \frac{18\phi'' + 8V+ 342V\phi''}
{4(1+V)\left(1+9\phi'' + 6V\!+54V\phi''\right)}.\label{eq:R_m-observe}
\end{equation}

However, $\mathcal{Y}_G, \mathcal{Y}_{G2}$ and $\mathcal{Y}_M$ linearly scale by the applied PSF differences ($\Sigma_0$ and $\Delta_0$). Nonetheless, if the secondary PSF is applied digitally to a single recording, the exact values of these differences are known. Therefore, we have a set of observables that are uniquely correlated to $V$, $\phi''$ and the digitally applied secondary PSF. We will show below how the PSF degree of freedom, allows us to choose a secondary PSF to optimize observable extraction \textit{vis-\`a-vis} the measurement noise. For the different sensitivity analyses and comparison to theory, we use alternatively the $\mathcal{J}$ and $\mathcal{Y}$ observables, having relative advantages per case, but being generally closely related.

The cumulant expansion method allows us to expand the application of the Guinier analysis, in which one fits the log-intensity versus $q^2$ curve to a polynomial form. Using cumulants, we derive the expression (see more detail in Eq.\,\eqref{eq:lnI_dimless_app} in Appendix \ref{app:cumulant}): 

\begin{equation}
{
\ln I(q)
= \ln I_0 - \frac{1}{3}\,(1+V)\,Q
\;+\;\left[\frac{2}{9}  V\;+\;\frac{1}{2}\phi''\,
\;+\;\,3\,V\phi''\right]\,Q^2
\;+\; O(Q^3).}
\label{eq:lnI_dimless_}
\end{equation}
where $Q=R_0^2 q^2$ (valid within the Guinier region, at $Q<1$). In the monodisperse case ($V=0$) the equation reduces to Eq.\,\eqref{eq:phi_series u}. By fitting the Guinier curve to a quadratic polynomial, $\ln I(q)\approx t_0+t_1q^2+t_2 q^4 $, we may now define a new observable 
\begin{equation}
    \mathcal{Y}_T\equiv \frac{t_2}{t_1^2} \approx\frac {9\phi''+4V +54V\phi''}{2(1+V)^2}\label{Y_T definition}
\end{equation}
This observable does not eliminate the instrumental PSF, but its functional form is independent of $R_0$. To first order in $V$, $\mathcal{Y}_T$ is not independent of $\mathcal{Y}_G$ (as $\mathcal{Y}_G-\mathcal{Y}_T\approx 1/4$). However, this observable is easy to obtain from the measurements and seems less susceptible to noise. The approach of extracting observables from experimental data and comparing them to results of forward-simulation for characterizing polydispersity resembles past works \cite{beaucage2004,tomchuk2023}, but requires fewer \textit{a priori} assumptions regarding the form-factor and the ensemble's size distribution shape.


\section{Numerical validation\label{sec:sim_valid}}
\subsection{Noise free simulation}
To validate our theory, we constructed simulations of the scattering from ensembles of particles. Each ensemble consisted of particles with a given monodisperse form-factor  and the statistics of the radius of gyration (mean and variance- $R_0$ and $V$). The simulation uses a realistic synchrotron beam setup (wavelength, distances, aperture sizes and detector resolution, see Tab.\ \ref{tab:beam_simul_param} in Appendix \ref{sensitivityAnalisys}). For each pixel we calculate the contribution of the scattering from each of the ensemble's members (according to its radius of gyration and the form-factor) and sum over all particles (abundance weighted). This produces a noise-free 2D intensity map of the scattering by a given ensemble. We perform a Guinier analysis and received the value of the apparent radius of gyration ($ \tilde{R}_0 $),
as discussed in Sec.\,\ref{sec:scattering-polydisp}. This is mainly needed for assessing the Guinier regime, and to perform the analysis only on pixels within it.

Next, we smeared the scattering maps by independently performing convolutions with anisotropic Gaussians (with different widths along the axes- $\sigma_x$ and $\sigma_y$). By comparing two such maps we obtained the log-ratio for each pixel, which may be represented in polar coordinates ($\mathcal{G}(q,\chi)$). The following step uses the least-squares method to fit $\mathcal{G}(q,\chi)$ to the radial and $\cos2\chi$ dependence (in the form shown in Eq.\,\eqref{eq:G_final}), thus obtaining the radial dependence of $G(q)$ and $M(q)$ 3\textsuperscript{rd}-order polynomials (e.\,g., see Fig.\,\ref{fig:2d m_g fit}). 
    
\begin{figure}
\begin{center}
\includegraphics[width=0.4
\paperwidth]{MGfitQcos2chi2.eps}
\end{center}

\caption{Simulation of the difference of log-intensity of a scattered intensity
after applying two different PSFs, plotted in azimuthal and radial axes ($\cos2\chi$
and $Q=q_{r}^{2}R_{g}^{2}$). The numerical fitting obtains a good
approximation and makes it easy to extract the functions $G$ and
$M$. \label{fig:2d m_g fit} The profile at $\cos2\chi=0$  corresponds to the functional form of $G(Q)$  in blue, while the red profiles at $\cos\chi=\pm1$  correspond to $G(Q)\pm M(Q)$.}
\end{figure}

Next, using the polynomial coefficient ratios of $G$ and $M$ and the Guinier analysis (Eqs.\,\eqref{eq:observables-old}-\eqref{eq:R_m-observe}), we can now calculate the observables  $\mathcal{J}_G(V,\phi'')$, $\mathcal{J}_{G2}(V,\phi'')$ and $\mathcal{J}_M(V,\phi'')$. 


For the validation of our analytical derivation we performed simulations over a range of $V\in[0,0.3]$ and either exact form-factors or using Guinier curvatures $\phi''\in[-0.02,0.06]$. Figure \ref{fig:theory-sim-disparity}  compares the calculation results to the theory, for the three PSF independent observables $\mathcal{J}$. For the sake of clarity, we present only the exact form-factors with extremal curvatures, those of Gaussian chains ($\phi''=1/18$) and spherical shells ($\phi''=-1/45$), along with one with an intermediate $\phi''=0.01$. The simulation of the $\mathcal{J}_G$ observable fits the theory in the three cases over a range of variances, reaffirming our theoretical derivation and the extraction method. 

However, the $\mathcal{J}_{G2}$ and $\mathcal{J}_M$ observables demonstrate a distinct disparity for the exact form-factors, and a good fit for the 'pure curvature' case up to a lower maximal variance. We explain this by the following argument: Exact form-factors have higher order curvatures ($\phi^{(3)}$, etc.) which were neglected in our theoretical derivation. The $\mathcal{J}_{G2}$ and $\mathcal{J}_M$  observables use higher order polynomial coefficients and therefore are more sensitive to the exact form-factors as well as being less accurate for higher $V$s. Nevertheless, Fig.\,\ref{fig:theory-sim-disparity} demonstrates that the actual observables' dependence on the variance is 'well-behaved' in all cases (including other form-factors that are not shown). This means that a single-valued calibration curve may be extracted if the form-factor is known. The $\mathcal{J}_{G}$ observable seems independent of the high order curvature, which makes it a more practical observable to use.


Now, let us discuss the dependence on other attributes of the ensemble. Previous works suggested polydispersity-related observables which showed a dependence on the variance as well as the distribution function \cite{tomchuk2023}. A prominent reason for the sensitivity to the distribution shape is that the scattering strength of a sphere is weighted by the sphere's mass squared, or $R_g^6$. Thus, a number distribution will have an observed scattering which fits a weighted mean and variance. In theory, in TENOR-SAXS the observables are independent of the distribution function, just to the weighted variance. We found that this holds quite well in practice. Specifically, we simulated scattering by ensembles of solid spheres, with different variances and different distribution functions (such as normal, lognormal, uniform, triangular, Boltzmann and Schulz distributions c.\,f.\, \cite{tomchuk2023}), and compared the calculated $\mathcal{Y}_G$ observable.

Figure \ref{fig:p_distribution_braid} shows how well the observable performs as an indicator for the weighted variance (right panels). For any given observation (with an unknown distribution function), the span of weighted variances of different distribution functions that produce the same observable value indicates the uncertainty of the weighted variance extraction, which is very small (compared to the sensitivity to other parameters, as will be shortly discussed). We note that the weighted variance in inherently smaller than the unweighted variance, since the weighting makes only a small range of $R_g$'s dominant. On the left panels, we can see the unweighted variance of the same simulations, showing a discrepancy which is inherent to SAXS, as adding small radii increases the unweighted variance, but their scattering strength is negligible.


\begin{figure}
\begin{center}
%\includegraphics[width=0.5\linewidth]{GvsV.eps}\includegraphics[width=0.5\linewidth]{MvsV.eps}
\includegraphics[width=\textwidth]{observ_theory.eps}

\end{center}
\caption{Observables obtained from 2D simulation results (dashed lines) and the theoretical first-order ratio from Tab.\,\ref{tab:GM-coef} (solid lines), plotted for a range of variances ($V$) and different form-factors. Observables of $\mathcal{J}_G$, $\mathcal{J}_{G2}$ and $\mathcal{J}_M$ are plotted in the left, center and right panels, respectively. In these simulations the normally distributed ensemble had an average $R_0=5\text{nm}$. in a typical Diamond B21 beamline setup, without noise. The theoretical curves are based on the Guinier curvature $\phi''$ only, so for the exact spherical shell and Gaussian chain form-factors who have higher orders the theory-simulation disparity is visible for the observables based on high-order polynomial fits- $\mathcal{J}_{G2}$ and $\mathcal{J}_M$.\label{fig:theory-sim-disparity}}
\end{figure}



The dependence of the observables on other parameters is discussed in detail in Appendix \ref{sensitivityAnalisys}. We tested the effect of different parameters within the approximation assumptions (Guinier $qR_g<1$ region, small enough instrumental PSF sizes, small variance). Importantly, we found that the instrumental PSF size hardly affects the observables, for different beam shapes and sizes, spanning between below typical synchrotron beam sizes to above laboratory X-ray setups. The PSFs' smearing sizes do change the observable values (mainly in high variances and large positive curvature), but since the smearing is done in a controlled way, this does not have an effect on the the obtained values. The observables are practically independent of the  and the average radius of gyration and are uniquely correlated to $V$ and $\phi''$. The ensemble's size (how many different particles are used for the simulation) is indistinct as well. 


\begin{figure}
\begin{center}
\includegraphics[width=1\textwidth]{V_braid_weighted_and_not_spherical_multi_dist_r4.eps}    

\end{center}
\caption{\label{fig:p_distribution_braid}The sensitivity to distribution type, in weighted and unweighted cases. Top left: Dependence of the $\mathcal{Y}_G$ observable on the polydispersity for different ensemble distribution functions (solid sphere's form-factor, $R_0=4nm$). The horizontal span of the curves' braid indicates the uncertainty of variance extraction on the distribution function (see annotated example). Bottom left: The span of $\sqrt{V}$ extraction as a function of the average $\sqrt{V}$. Right panels (top and bottom) present the same data scaled by the scattering strength weighted normalized variance, $V$ ($\propto R_g^6$).}
\end{figure}


\subsection{Noisy simulation analysis}
Having produced the direct dependence of the ensemble properties on observables produced from the log-ratio of measured intensities with two anisotropic PSFs, we may sketch a recipe for using the method for practical application.
In order to obtain two different PSFs, one may use several methods, such as repeating the measurement with different slit sizes or different sample-detector distances. However, a simpler route would be to take a 2D measurement and apply a digital filter, as we used in the simulations. This method is more practical for several reasons. Notably, it eliminates the need for multiple costly experiments, reduces concerns about sample stability, and avoids dependencies on fixed beam centers, slit settings, repeated buffer subtractions, normalization or other system parameters that may vary between measurements.

Real measurements are noisy and depend on different parameters such as measurement time, source and sample stability, and more. Therefore, we do not know \textit{a priori} which smearing would be best for each observable fit. However, it is easy to grade the polynomial $G$ and $M$ (Eq.\,\eqref{eq:GM-poly} fits' quality of different PSF sets by comparing the standard errors for the observable (Eqs.\,\eqref{eq:observables}-\eqref{eq:R_m-observe}). Thus, by scanning the smearing sizes, it is possible to find the best estimate of the observable available from the data.\\

\begin{comment}
\begin{figure}\begin{center}
\includegraphics[width=0.45\textwidth]{norm_R_g3_V30Nu_02.eps}   
\includegraphics[width=0.45\textwidth]{norm_R_g5_V30Nu_02.eps}   
\includegraphics[width=0.45\textwidth]{Normal_R_g5_V30Nu06.eps}   
\includegraphics[width=0.45\textwidth]{Normal_R_g5_V00Nu06.eps}   
\includegraphics[width=0.45\textwidth]{uniform_R_g5_V30Nu_02.eps}   
\includegraphics[width=0.45\textwidth]{triang_R_g5_V30Nu_02.eps}   
\end{center}
\caption{\label{fig:V-intersect}Curves of the estimated $V$ as a function of the assumed curvature $\phi''$ for the observables, with negligible noise and Diamond B21-like PSF. The intersections allow extraction of both $V$ and $\phi''$. The marker indicates the actual values used in each example. The observables were calculated using a variety of radius distribution functions \cite{tomchuk2023}, while the ground truth used for creation of the plots used normal distribution and an ideal PSF. Nevertheless, the intersection $V$  values differ by about $0.01$ from the actual $V$. The examples include different average $R_g$, and different variances and curvatures (typically $V=0.3$ and $\phi''=-0.02$ have the largest curve slopes).}
\end{figure}
\end{comment}

The next step is extraction of the model parameters using simulations of ground truths of different variances. Repeating the process with simulated data (keeping the same $R_g \approx \tilde{R}_0 $ as the one obtained by the Guinier analysis, i.\,e., setting the average radius to be $\tilde{R}_0/\sqrt{1+V}$ for the selected $V$), for the chosen set of digitally applied PSFs, we can find the observable functions for noise-free ground truths, with the same form-factor of the particles: $\mathcal{Y}_G^{GT}(V)$ (the scheme may be executed also for the other observables, $\mathcal{Y}_{G2}$ and $\mathcal{Y}_M$, but with lesser accuracy). Moreover, we can find the reciprocal relation between the variance and the observable $V^{GT}_{G}$ ($\mathcal{Y}^{GT}_{G}$). This is effectively a calibration function, which allows direct extraction of the ensemble's polydispersity from the observable value. %$\mathcal{Y}_{G/G2/M}^{GT}(V,\phi'')=\mathcal{Y}_{G/G2/M}$. 
Extraction of the form-factor from the combination of different observables is possible in theory, but in practice requires extremely low noise levels.
%Practically, for an initial rough estimation of the form-factor curvature ($\phi''$), we look for intersections between obtained curves of $V^{GT}_{G/G2/M}(\phi'')$ using the values of measured $\mathcal{Y}_{G/G2/M}$ (Fig.\,\ref{fig:V-intersect}, note that the extrapolation may bring about negative values of $V$). Not always there are intersections, since (a) noise may sway the observables and the estimated $R_g$, and (b) the real measurement PSF provides different results than the ideal PSF used for the simulation. 
\section{TENOR-SAXS analysis recipe and results}

Solution SAXS data is typically conducted on the azimuthal 1D integrated datasets, neglecting the PSF contribution and sample variability. Here, we suggest a recipe to extract $R_g$ and $V$ knowing the underline form factor. Uniquely, TENOR-SAXS analysis is unbiased to the PSF contribution, and performs well at feasible photon counts. Moreover, it is sensitivity to different particle distribution function is small. Below we list the key steps to perform TENOR-SAXS analysis.
\begin{enumerate}
\item Obtain a 2D SAXS measurement. After removing the background intensity, the data should have the intensity for the 2D scattering angle $I(\mathbf{q})$. Note that background measurement removal is typically done in radially averaged (1D) data, and here the 2D process should be done with care, for example, by projecting the averaged 1D background data back to 2D and then subtracting it.  
\item Perform a Guinier analysis of the data, thus obtaining $\tilde{R}_0\approx R_0 \sqrt{(1+V)}$ (assuming that the original PSF has been small enough not to affect this value). 
\item  Scan PSF sets to find best fitted observable values:
\begin{enumerate}

\item Scan over PSFs: $\sigma_{x_{1,2}}$ and $\sigma_{y_{1,2}}$ which are within the Guinier region ($\sigma \tilde{R}_0\ll1$). Create a (typically Gaussian) kernel for each PSF. In practice, two pairs of pixel indices (a quartet) are chosen, and each pair specifies the dimensions of an independent 2D Gaussian kernel, yielding two distinct kernels. Typically, the best smearing kernel may be predetermined, tending toward large areas (over which the integration reduces noise), but low enough so that there are enough pixels left that are not affected by data from outside the Guinier region. As described above, the difference between the two PSFs is important as well, so typically the size of second PSF is about half of the first one's. In the predetermined case, one may limit the scan to just a few cases in the vicinity of the chosen quartet.
\item Apply a digital filter- convolve with the kernel obtaining two intensity maps, differing by their PSFs.
\item Take the logarithm of the ratio of the two maps, and perform a weighted fit the resulting $\mathcal{G}(q,\chi)$ to the form of Eq.\,\eqref{eq:G_final} and to the 2\textsuperscript{nd} order polynomials $G(q)$ and $M(q)$. The weight of each pixel is the square root of its intensity.
\item Choose the PSF quartet that leads to best fit for each observable using Eqs.\,\\\eqref{eq:observables}-\eqref{eq:R_m-observe}, extacting $\tilde{R}_0$, and the polynomials coefficients of $G(q)$ and $M(q)$, from which the observables $\mathcal{Y}$ are obtained.
\end{enumerate}
\item Find $V$ for each observable $\mathcal{Y}$:
\begin{enumerate}
    \item Create the ground truth- simulate the scattered intensity of an ensemble with the known form-factor (an explicit function or using the curvature $\phi''$), the measured $R_g$  and different $V$ values. Convolve with the same kernels obtained above ('the best-fit quartet'). Extract the observable as done above.
    \item Interpolate to find the $V$ that reproduces the value of $\mathcal{Y}$.
\end{enumerate}
%\item Each intersection of two $V(\phi'')$ curves provides a value of $V$ (and $\phi''$). Take the average of the intersections' $V$ values as the estimate.


\end{enumerate}

The above recipe requires the following assumptions to be followed: 
\begin{enumerate}
    \item The $q$-range should be in the Guinier region. This may be verified with the Guinier analysis result, limiting the 2D analysis to $q^2\cdot R_0^2(1+V) < 1$. Note that it was found that sometimes (high variances or form-factors with large third derivative $\phi'''$) it was found best to limit the analysis to about 90\% of that $q$-range.
    \item The PSF width should be small ($ \sigma^2\cdot R_0^2(1+V)\ll1 $), where \(\sigma\) refers to the scattering vector smearing, typically due to beam divergence and angular uncertainty $\Delta \theta$ related to aperture and pixel sizes,  $\sigma\approx 4\pi/\lambda\cdot \Delta\theta$.
    \item The PSF anisotropy should be different for the two 2D maps, which is generally true if we compare two measurements (the system including the beam profile and slit size is never totally isotropic). In the digital filtering option, one should typically use a distinctly anisotropic filter. Generally, increasing the anisotropy difference provides better results.  
    \item The performance changes with PSF size, and widening the PSF introduces an interplay between the effect of noise reduction due to averaging over a large area and the reduction of the number of independent pixels for the 2D fitting. Fortunately, one may quantify said fitting quality, and thus, by repeating the process with different PSF combinations (each producing a different result), we know which combination is the best, and use it's result to extract the relative variance. 
\end{enumerate}
We performed numerous simulations and scanned the applicable regime (in the Diamond B21 beamline instrument configuration, c.f.\,Tab.\,\ref{tab:beam_simul_param} in App.\,\ref{sensitivityAnalisys}) spanning $-0.02<\phi''<+0.06$,  $0<V<0.3$, $2~\text{nm} < R_0 < 8~\text{nm}$ and  with Gaussian (normal) probability distribution function. Applying the TENOR-SAXS method provided the sensitivity to different simulation and ensemble parameters, as well as to the TENOR-SAXS' specifications, showing that the ground truth (noise-free) observables are mostly insensitive to them, details are presented in App. \ref{sensitivityAnalisys}. Noteworthy is the result that lowering $R_0$ mainly affects the results through increasing the number of pixels within the Guinier regime, thus the 2D fit is based on more points, reducing the fit error. Additionally, more photons are accumulated withing the Guinier region, increasing the signal-to-noise ratio.

We have performed the TENOR-SAXS recipe with simulated data with different noise levels. The noise model was Gaussian with a standard deviation equal to the square root of the number of photons per pixel \cite{gutman2020}. For this simulation we used a full solid-sphere's form-factor \cite{guinier1955}, with log-normal radius distribution, as is typical for aerosol particulates \cite{LANDGREBE1990,beaucage2004}. The scattering intensity was weighted by $R_g^6$. As expected, for different numbers of photons (and hence noise level) we see an increase in performance as the noise level decreases. We estimated the variance \(V_{est}\) based on the $\mathcal{Y}_G$ observable, by comparing to a calibration using the observable obtained from the ground truth simulation. Figure \ref{fig:V-discrep} shows the histograms of deviation between the input and the estimated $\sqrt{V}$ (distribution relative standard deviation) for different forward scattering photons. The results show that when the peak collected forward scattered photon density in the $q$-space is more than $3\times 10^8\,\mathrm{photons/nm^{-2}}$ over 90\%  of the measurements show valid variances (the observable is within the calibration range). The estimated distribution spread (with 5\% outlier omission) seems to be unbiased, with a discrepancy standard deviation lower than 0.1 in $\sqrt{V}$. As detailed in App.\,\ref{sensitivityAnalisys}, the photons needed for this performance are collected in an order of one-to-few frames (seconds of measurement) in the Diamond light source BioSAXS B21 beamline with a regular sample, making TENOR-SAXS a feasible method.

\begin{figure}
%\includegraphics[width=1.1\textwidth]{Yg_discrep_3nois.eps}
\includegraphics[width=1.1\textwidth]{lognormal_spherical_rg4_violin_plot.eps}
\caption{\label{fig:V-discrep}
The discrepancy statistics of simulations of ensembles of solid sphere form-factor. The ensemble's radius of gyration distribution is log-normal with a mean of $R_0=4\,\mathrm{nm}$ and relative standard deviations in the range of 0-0.55. The $\mathcal{Y}_G$ obtained from each simulation were compared with that obtained from an ideal (noiseless and point-like aperture) ground truth simulation. In this violin plot the patches represent the histogram of the discrepancy of the derived relative standard deviation ($\sqrt{V}$) from the actual value for each signal strength. 5\% outliers were omitted from the analysis. The mean bias is represented by a solid line, while the discrepancy standard deviation margin is indicated by the dashed line. The percentage of valid points (where the observable is not out of the valid range of the ground-truth observable, indication the method failed) is indicated on each histogram.
%The discrepancy (over a range of ensembles and scattering curvatures- $V\in[0,0.3]$ and $\phi''\in[-0.02,0.06]$) of the detected size distribution's relative standard deviation $p$ compared to the simulated ground truth, for different noise levels (collected photon number densities). The solid line indicates the average discrepancy, while the shaded area indicates the $1\sigma$  region around it. The estimation is based on the $\mathcal{Y}_G$ observable when $\phi''$ is given. The simulation is based on Diamond's B21 BioSAXS beam's published parameters (where the scattered photon density may reach  $10^9$ photons/$\mathrm{nm}^{-2}$, with a sample of $R_0=5nm$.
}
\end{figure}

\section{Summary and Conclusions}

We presented here a Technique for ENsemble Observation by Resolution-variation in SAXS (TENOR-SAXS). It  provides a practical means to quantify ensemble variance directly from two-dimensional scattering images while removing explicit dependence on instrumental resolution functions. By exploiting digital smearing, the method retrieves distributional information in polydisperse samples that is not accessible through conventional Guinier analysis. In contrast to ensemble-modeling strategies that fit large structural pools to one-dimensional profiles~\cite{bernado2007,tria2015,antonov2016}, or density-reconstruction approaches that recover three-dimensional real-space maps~\cite{grant2018}, TENOR-SAXS employs PSF-independent observables obtained from controlled resolution variation. The technique thus enables ensemble characterization without imposing heavy structural priors, complementing both model-based and model-free SAXS methodologies~\cite{hura2009,rambo2013,burger2016}.

TENOR-SAXS offers two principal advantages: practicality and robustness. Digital smearing eliminates the need for repeated measurements under modified beamline conditions, thereby reducing beamtime requirements, alignment and normalization errors. Moreover, in some cases the technique requires no \textit{a priori} knowledge on the exmained system such as the probability distribution function shape or the underlined sample geometry. Our results indicate that systematic deviations remain negligible under realistic conditions, with ensemble relative standard deviation recovered within $\Delta \sqrt{V}  < 0.1$ when scattered-photon densities exceeding $N_{q^2} \sim 3\cdot10^{8}$~photons/nm$^{-2}$. Such levels are readily achievable at modern synchrotron facilities and accessible for advanced laboratory SAXS instruments. Because the analysis operates directly on raw 2D intensity maps, the method may also be applied retrospectively to archived datasets, broadening the interpretative value of existing data.

The approach has some limitations. Its current formulation applies to the Guinier regime, requires modest anisotropic differences in smearing, and depends on the curvature term in the Guinier expansion ($\phi''$). Noise sensitivity increases at low photon counts, sometimes preventing the variance to be extracted from the measurement. However, the noise dependency is predictable and allows experiments to be designed accordingly. Buffer subtraction, typically performed after azimuthal averaging, must be carried out consistently at the two-dimensional level to avoid introducing anisotropic artifacts in the log-ratio observables.

Several avenues for further development are apparent. Alternative observables may reduce dependence on $\phi''$ and enhance noise suppression. Extension to higher $q$-ranges and more complex form factors would broaden applicability. Systematic mapping of dependencies on point-spread function size, detector resolution, and fitting strategies may provide practical guidelines for experiment design. Incorporating independent structural information, such as externally determined $\phi''$ or orthogonal fitted parameters, may further reduce uncertainty. Further analysis is needed for partical conformations where the form-factor curvature is related to the radius of gyration. Most importantly, experimental validation is now required. Application to new measurements or re-analysis of archived 2D images is called for, with particular emphasis on systems where structural heterogeneity plays a functional role, such as intrinsically disordered proteins and polydisperse nanomaterials.

In summary, TENOR-SAXS introduces a conceptually simple, PSF-independent framework for ensemble characterization in SAXS. By combining practical implementation, robustness to systematic effects, and complementarity with existing approaches, it has the potential to substantially expand the interpretative power of small-angle scattering in both biological and materials sciences.

\section*{Methods}
 Methods and any associated references are available in the online version of the paper.

\section*{Acknowledgments}
This research was supported by the Pazy Research Foundation (341- 2022).  

\section*{Appendix}
\appendix


\section{Direct derivation of the second order log-intensity}
\label{app:direct-logI}

We start from the definition of the measured intensity as a convolution of the intrinsic form factor with the instrumental PSF:
\begin{equation}
I(\mathbf q) \;=\;  R (q) \circledast F(q)  \;=\; \int R(-\mathbf p)\,F(\mathbf q+\mathbf p)\,d^2p,
\qquad
\int R(\mathbf p)\,d^2p=1.
\label{eq:conv-direct}
\end{equation}
Assume \(R\) is a centered, separable anisotropic Gaussian,
\begin{equation}
\begin{aligned}
R(\mathbf q) \;=\; \frac{1}{2\pi\sigma_x\sigma_y}\,
\exp\!\Big(-\frac{q_x^2}{2\sigma_x^2}-\frac{q_y^2}{2\sigma_y^2}\Big),
\qquad\\
\langle q_x\rangle=\langle p_y\rangle=0,\;\;
\langle q_x^2\rangle=\sigma_x^2,\;\;\langle q_y^2\rangle=\sigma_y^2,\;\;\langle q_xq_y\rangle=0.
\label{eq:R-gauss-direct}
\end{aligned}
\end{equation}
%\subsection*{Step 1: moment expansion of the convolution}
Expand \(F(\mathbf q+\mathbf p)\) to second order:
\begin{equation}
F(\mathbf q+\mathbf p) \;\approx\; F(\mathbf q)
+ p_i\,\partial_{q_i}F(\mathbf q)
+ \frac{1}{2}\,p_ip_j\,\partial_{q_i}\partial_{q_j}F(\mathbf q),
\qquad i,j\in\{x,y\}.
\label{eq:taylor-F}
\end{equation}
Averaging against the even kernel \(R\) eliminates the linear term. Using the second moments in Eq.~\eqref{eq:R-gauss-direct} gives
\begin{equation}
I(\mathbf q) \;\approx\; F(\mathbf q) \;+\; \frac{1}{2}\Big[\sigma_x^2\,\partial_{q_x}^2+\sigma_y^2\,\partial_{q_y}^2\Big]F(\mathbf q).
\label{eq:I-second-order}
\end{equation}
This is the standard small-PSF result. More generally, for any centered kernel with covariance matrix \(C\), one has \(I \approx F + \tfrac{1}{2}C_{ij}\partial_{q_i}\partial_{q_j}F\).

%\subsection{Step 2: taking the logarithm with a dimensionless small parameter}
Next, we write Eq.~\eqref{eq:I-second-order} as \(I=F+\delta\) with
\begin{equation}
\frac{\delta}{F} \;=\; \frac{1}{2}\left[\sigma_x^2\,\frac{\partial_{q_x}^2F}{F}+\sigma_y^2\,\frac{\partial_{q_y}^2F}{F}\right].
\label{eq:delta-over-F}
\end{equation}
Following, we define the dimensionless smallness measure as:
\begin{equation}
\varepsilon(\mathbf q) \;\equiv\; \max\Big\{
\big|\sigma_x^2\,\partial_{q_x}^2\ln F(\mathbf q)\big|,\;
\big|\sigma_y^2\,\partial_{q_y}^2\ln F(\mathbf q)\big|,\;
\big|\sigma_x^2\big(\partial_{q_x}\ln F(\mathbf q)\big)^2\big|,\;
\big|\sigma_y^2\big(\partial_{q_y}\ln F(\mathbf q)\big)^2\big|
\Big\} \;\ll\; 1,
\end{equation}
so that \(\delta/F = \mathcal O\!\big(\varepsilon(\mathbf q)\big)\) is unitless.  Since \(\ln(1+x)=x+\mathcal O(x^2)\), the expanded log scattered intensity becomes: 
\begin{equation}
\ln I(\mathbf q) \;=\; \ln F(\mathbf q) + \ln\!\Big(1+\frac{\delta}{F}\Big)
\;\approx\; \ln F(\mathbf q) + \frac{\delta}{F}
\;+\; \mathcal O\!\big(\varepsilon(\mathbf q)^2\big).
\label{eq:logI-expansion}
\end{equation}
Using the identity
\begin{equation}
\frac{1}{F}\,\partial_{q_i}^2F \;=\; \partial_{q_i}^2\ln F \;+\; \big(\partial_{q_i}\ln F\big)^2,
\qquad i\in\{x,y\},
\end{equation}
and substituting into Eq.\,\eqref{eq:logI-expansion} yields the second order log-intensity formula in a fully dimensionless form:

\begin{equation}
\begin{aligned}
\ln I(\mathbf q)\;\approx\;&\;\ln F(q) + \tfrac{1}{2}\big(\sigma_x^2\,\partial_{q_x}^2+\sigma_y^2\,\partial_{q_y}^2\big)\ln F(q) \\[6pt]
&+ \tfrac{1}{2}\Big[\sigma_x^2\big(\partial_{q_x}\ln F(q)\big)^2
     + \sigma_y^2\big(\partial_{q_y}\ln F(q)\big)^2\Big] + \mathcal O\!\big(\varepsilon(\mathbf q)^2\big)
\end{aligned}%
\label{eq:logI-operator-appendix-updated}
\end{equation}

The neglected terms are of second order in the dimensionless combinations \(\sigma_i^2\,\partial_{q_i}^2\ln F\) and \(\sigma_i^2\,(\partial_{q_i}\ln F)^2\). This is met when the angular smearing is smaller than the form-factor's second log-$q$-derivative which is typically comparable to $R_g$. That is, the PSF requirement for this approximation is $2\sigma_i^2 R_0 ^2/3 \ll 1$ at the Guinier region.

For an isotropic form factor \(F(\mathbf q)=F(| \mathbf q|)\), define the first and second radial derivatives
\[
f(q)\equiv \partial_q\ln F(q),\qquad f'(q)\equiv \partial_q^2\ln F(q).
\]

Projecting the Cartesian derivatives onto the radial direction then yields
\begin{align}
\partial_{q_x}\ln F &= f(q)\frac{q_x}{q},\quad
\partial_{q_y}\ln F = f(q)\frac{q_y}{q},\\
\partial_{q_x}^2\ln F &= f'(q)\frac{q_x^2}{q^2}+\frac{f(q)}{q}\Big(1-\frac{q_x^2}{q^2}\Big),\quad
\partial_{q_y}^2\ln F = f'(q)\frac{q_y^2}{q^2}+\frac{f(q)}{q}\Big(1-\frac{q_y^2}{q^2}\Big).
\end{align}

Next, we define the azimuth $\chi$ by $(\frac{q_x}{q}, \frac{q_y}{q})=(\cos\chi,\sin\chi)$ and the effective radial smearing at each point by: 
\begin{align}
\Sigma_{\mathrm{eff},k}^2(\chi)
&\equiv\sigma_{x,k}^2\cos^2\chi+\sigma_{y,k}^2\sin^2\chi .
\label{eq:sigma_1}
\end{align}
Following, we can now transform Eq.\,\eqref{eq:logI-operator-appendix-updated} to obtain the Cartesian second order result:
\begin{equation}
\ln I_k(q,\chi)\approx \ln F(q)+\frac{1}{2}\Big[
\Big(f'(q)+f(q)^2\Big)\,\Sigma_{\mathrm{eff},k}^2(\chi)
+\frac{f(q)}{q}\Big(\sigma_{x,k}^2+\sigma_{y,k}^2-\Sigma_{\mathrm{eff},k}^2(\chi)\Big)
\Big].
\label{eq:PSF-cart-core}
\end{equation}

Next, from geometrical considerations ($\cos^2\chi=(1+\cos 2\chi)/2$ and $\sin^2\chi=(1-\cos 2\chi)/2$), we can rewrite Eq.\,\eqref{eq:sigma_1}:
\begin{equation}
\Sigma_{\mathrm{eff},k}^2(\chi)=\bar\sigma_k^2+\frac{\Delta\sigma_k^2}{2}\cos 2\chi,\quad
\sigma_{x,k}^2+\sigma_{y,k}^2-\Sigma_{\mathrm{eff},k}^2(\chi)=\bar\sigma_k^2-\frac{\Delta\sigma_k^2}{2}\cos 2\chi,
\end{equation}
where
\begin{equation}
\bar\sigma_k^2=\frac{\sigma_{x,k}^2+\sigma_{y,k}^2}{2},\qquad \Delta\sigma_k^2=\sigma_{x,k}^2-\sigma_{y,k}^2.
\end{equation}
Next, we define the PSF scalars:
\begin{equation}
\Sigma_{0,k}\equiv \bar\sigma_k^2,\qquad \Delta_{0,k}\equiv \frac{\Delta\sigma_k^2}{2},
\label{eq:Sigma0-Delta0}
\end{equation}
and the two radial response functions:
\begin{equation}
G(q)\equiv \frac{f(q)}{q}+f'(q)+f(q)^2,\qquad
M(q)\equiv -\frac{f(q)}{q}+f'(q)+f(q)^2.
\label{eq:G-M-defs}
\end{equation}
Then Eq.\,\eqref{eq:PSF-cart-core} becomes the pixel-wise model
\begin{equation}
I_k(q,\chi)\approx F(q)+\mathcal{G}_k(q,\chi),\quad
\mathcal{G}_k(q,\chi)=\frac{1}{2}\,\Sigma_{0,k}\,G(q)+\frac{1}{2}\,\Delta_{0,k}\,M(q)\cos 2\chi
\label{eq:logIk}
\end{equation}
which is the form used for the 2D parameter extraction. Specifically, for two measurements with different PSFs,
subtracting Eq.\,\eqref{eq:logIk} eliminates $\ln F$ and yields:
\begin{equation}
\big[\ln I\big]_1-\big[\ln I\big]_2
=\frac{1}{2}\Big(\Sigma_{0,1}-\Sigma_{0,2}\Big)\,G(q)
+\frac{1}{2}\Big(\Delta_{0,1}-\Delta_{0,2}\Big)\,M(q)\cos 2\chi,
\label{eq:G-cart}
\end{equation}
which separates the angle independent and $\cos 2\chi$ components and depends only on the differences of the PSF scalars and the common radial functions $G$ and $M$.

%roy is here

\[
\phi''_{\textbf{eff}}{(0)}=\beta+\alpha, \qquad
\]
\[
\phi''_{\text{eff}}(V)=3\beta+5\alpha
\]
\section{Cumulant expansion of the ensemble-averaged intensity}
\label{app:cumulant}

\subsection{Notation and dimensionless variables}
We consider isotropic single-particle scattering, so the intrinsic (monodisperse) form factor depends on the radial scattering vector only via $q^2$. For an ensemble with a distribution of radii of gyration $R_g$, we define the random variable $X \equiv R_g^2$, and the dimensionless argument $u \equiv q^2 X$. \\

Let $\phi(u)$ denote the \emph{monodisperse} log-form factor,
\[
\ln F(q\,|\,X) \equiv \phi(u),
\]
assumed analytic at $u=0$. Its Taylor coefficients at the origin are
\begin{equation}
\phi(u)= \kappa\,u + \tfrac{1}{2}\phi''\,u^2 + \tfrac{1}{6}\phi^{(3)}\,u^3 + O(u^4),
\qquad
\kappa=\partial\phi(0),\;\; \phi''=\partial^2\phi(0),\;\; \phi^{(3)}=\partial^3\phi(0),
\label{eq:phi_series_app}
\end{equation}
and relative to the intensity at $q=0$, the ensemble-averaged measured intensity is: 
\begin{equation}
I(q)=\mathbb{E}_X\!\Big[\exp\!\big(\phi(q^2 X)\big)\Big] = \mathbb{E}_X\!\big[e^{\,\phi(u)}\big], 
\qquad
\ln I(q)=\ln \mathbb{E}_X\!\big[e^{\,\phi(u)}\big].
\label{eq:def_lnI_app}
\end{equation}
\subsection{Cumulant-generating function}
We introduce the random variable $Z \equiv \phi(u)$ and its cumulant-generating function (CGF):
\begin{equation}
K_Z(s)=\ln \mathbb{E}[e^{s Z}].    
\end{equation}
Following, $\ln I(q)=K_Z(1)$ and therefore we can expand the logarithm such that: 
\begin{equation}
\ln I(q)=\sum_{n\ge 1}\frac{\kappa_n(Z)}{n!}
=\underbrace{\mathbb{E}[Z]}_{\kappa_1(Z)}+\frac{1}{2}\underbrace{\mathrm{Var}(Z)}_{\kappa_2(Z)}
+\frac{1}{6}\kappa_3(Z)+\mathcal{O}(u^4).
\label{eq:cum_series_app}
\end{equation}
Using Eq. \eqref{eq:phi_series_app} up to $\mathcal{O}(u^3)$ we can rewrite our random variable as: 
\begin{equation}
Z = A\,u + B\,u^2 + C\,u^3 + \mathcal{O}(u^4),
\qquad
A=\kappa,\quad B=\tfrac{1}{2}\phi'',\quad C=\tfrac{1}{6}\phi^{(3)}.
\label{eq:Z_ABC_app}
\end{equation}
%\paragraph{First cumulant.}
According to cumulant expansion theorem, we can express the cumulants using Eq. \ref{eq:Z_ABC_app}. The first cumulant is therefore:  
\begin{equation}
\kappa_1(Z)=\mathbb{E}[Z]
= A\,\mathbb{E}[u] + B\,\mathbb{E}[u^2] + C\,\mathbb{E}[u^3] + \mathcal{O}(u^4).
\label{eq:k1_app}
\end{equation}
To $\mathcal{O}(u^3)$ one needs only terms up to $u^3$ in $Z$ and bilinears up to $u^3$ to find the second cumulant:
\begin{equation}
\kappa_2(Z)=\mathrm{Var}(Z)=A^2\,\mathrm{Var}(u)+2AB\,\mathrm{Cov}(u,u^2) + O(u^4).
\label{eq:varZ_app}
\end{equation}
And last, to a leading term of $\mathcal{O}(u^3)$, the third cumulant does \emph{not} require $C\neq 0$ and can be expressed as:
\begin{equation}
\kappa_3(Z)=\mathbb{E}\!\big[(Z-\mathbb{E}[Z])^3\big]
=(A u - A\,\mathbb{E}[u])^3 + \mathcal{O}(u^4)
= A^3\,\kappa_3(u) + \mathcal{O}(u^4).
\end{equation}
For a given $q$, the random variable is $X$, and therefore $\kappa_3(u)=q^6\,\kappa_3(X)$. For a narrow, nearly symmetric size distribution, $\kappa_3(X)=\mathcal{O}(V^{3/2})$, hence
\begin{equation}
\kappa_3(Z) = \mathcal{O}\!\big(u^3\,V^{3/2}\big).
\label{eq:k3_drop_app}
\end{equation}
Thus, even if $\phi^{(3)}=0$ (\textit{i.e.,} $C=0$), the $\mathcal{O}(u^3)$ contribution survives. However, for  small values of $V$, this component originating from distribution skewness can be neglected.
\\
Next, we evaluate the leading moments in terms of $X$ for a given $q$:
\begin{equation}
\begin{aligned}
\mathbb{E}[u]=q^2\,\mu_1,\quad \mathbb{E}[u^2]=q^4\,\mu_2,\quad \mathbb{E}[u^3]=q^6\,\mu_3, 
\qquad \\
\mathrm{Var}(u)=q^4\,\mathrm{Var}(X),\quad \mathrm{Cov}(u,u^2)=q^6\,\mathrm{Cov}(X,X^2),
\label{eq:u_moments_app}
\end{aligned}
\end{equation}
where $\mu_k\equiv \mathbb{E}[X^k]$. For a narrow ensemble of $R_g$ with mean $R_0=\mathbb{E}[R_g]$ and dimensionless variance $V=\mathrm{Var}(R_g)/R_0^2\ll 1$, the required combinations to first order in $V$ are
\begin{equation}
\mu_1=R_0^2(1+V),\qquad 
\mu_2\approx R_0^4(1+6V),
\label{eq:X_moments_app}
\end{equation}
\begin{equation}
\mathrm{Var}(X)\approx 4R_0^4 V,\qquad 
\mathrm{Cov}(X,X^2)\approx 8R_0^6 V.
\label{eq:cov_app}
\end{equation}

\subsection{Result for $\ln I(q)$ through $\mathcal{O}(u^3)$}
Combining \eqref{eq:cum_series_app}-\eqref{eq:cov_app} and keeping terms through $\mathcal{O}(u^3)$ and linear in $V$,
\begin{align}
\ln I(q)
&= A\,\mathbb{E}[u] + B\,\mathbb{E}[u^2] + \frac{1}{2}A^2\,\mathrm{Var}(u) + A B\,\mathrm{Cov}(u,u^2) 
+ \mathcal{O}(u^4) \nonumber\\[2pt]
&= \kappa\,q^2\,\mu_1 
\;+\; \tfrac{1}{2}\phi''\,q^4\,\mu_2
\;+\; \tfrac{1}{2}\kappa^2\,q^4\,\mathrm{Var}(X)
\;+\; \kappa\,\tfrac{1}{2}\phi''\,q^6\,\mathrm{Cov}(X,X^2)
\;+\; \mathcal{O}(u^4).
\label{eq:lnI_u_app}
\end{align}
Using \eqref{eq:X_moments_app}-\eqref{eq:cov_app} and defining $Q\equiv q^2 R_0^2$,
\begin{equation}
\ln I(q)
= \kappa\,(1+V)\,Q
\;+\;\left[\tfrac{1}{2}\phi''\,(1+6V)
\;+\;2\kappa^2  V\,\right]\,Q^2
\;+\;4\kappa\phi''  V\,Q^3
\;+\; O(Q^4), \\
\label{eq:lnI_dimless_app}
\end{equation}
which limits the analysis to sub-Guinier region ($qR_0<1$). We can further rewrite Eq. \ref{eq:lnI_dimless_app} as a polynomial in $Q$: 
\begin{equation}
\ln I(q) \;=\; a\,Q \;+\; b\,Q^2 \;+\; c\,Q^3 \;+\; \mathcal{O}(Q^4),
\label{eq:abc_Q_app}
\end{equation}
where, to first order in $V$:

\begin{align}
a &= \kappa\,(1+V), \\
b &= \tfrac{1}{2}\phi''\,(1+6V) + 2\kappa^2\,V, \\
c &= 4\kappa\phi'' \, V. 
\end{align}



\subsection{From cumulant expansion to $G$ and $M$ functions}

We are now in position to express the polynomial functions $G(q)$ and $M(q)$ in Eq. \ref{eq:GM-poly} using the cumulant expansion in Eq. \ref{eq:abc_Q_app}. 
In the main text, we identified the polynomial function (Eq. \ref{eq:GM_defs}) as 
\(
G(q)=f'(q)+\tfrac{f(q)}{q}+f(q)^2
\)
and 
\(
M(q)=f'(q)-\tfrac{f(q)}{q}+f(q)^2
\),
where $f(q)=\partial_{q}\ln I(q)$. Using Eq. \eqref{eq:abc_Q_app} and $Q=q^2 R_0^2$, the expansions in powers of $q^2$ are
\begin{align}
G(q) &= 4a\,R_0^2 \;+\; (16b+4a^2)\,R_0^4\,q^2 \;+\; (36c+16ab)\,R_0^6\,q^4
\;+\; (24ac+16b^2)\,R_0^8\,q^6 \;+\; \mathcal{O}(q^8), \label{eq:G_from_abc_app}\\
M(q) &= 0 \;+\; (8b+4a^2)\,R_0^4\,q^2 \;+\; (24c+16ab)\,R_0^6\,q^4
\;+\; (24ac+16b^2)\,R_0^8\,q^6 \;+\; \mathcal{O}(q^8). \label{eq:M_from_abc_app}
\end{align}
Equations \eqref{eq:lnI_dimless_app}-\eqref{eq:M_from_abc_app} provide the dimensionless, order-by-order link between the ensemble cumulants (via $a,b,c$) and the measured 2D coefficients.

Keeping terms linear in $V$, and assuming  $\phi^{(3)}\approx0$ (as we are interested in the low $q$ regime) we can now use the cumulant expansion to identify the relation between $a,b,c$ and the polynomial coefficients for $G(q)$ and $M(q)$ shown in Table \ref{tab:GM-coef} where we show expressions truncated at $\mathcal{O}(V)$ and $\mathcal{O}(q^6)$. We note that in some cases, for example where $\phi '' = 0$, the expansion requires to take into account $\phi^{(3)}$ contribution. In general the results presented in Table \ref{tab:GM-coef} are a good first approximation to many cases, but appropriate cumulant expansion should be practiced for the relevant form-factor under investigation. In practice, given a relevant form-factor, one can simulated the ideal scattering for various $V$ to find the relevant coefficients $G(q,V)$ and $M(q,V)$. This can be used to produce a calibration curve from the the observables (Eqs. \ref{eq:observables}-\ref{eq:R_m-observe}) to $V$.          

\section{Degree of independence on ensemble and instrumental parameters}
\label{sensitivityAnalisys}In order to check how the observables depend on different parameters, and to find the suitability of the method to be compared to a simulation (calibration), we simulated the scattering for different ensembles, and extracted the observables. We varied either the average radius of gyration, the number of different ensemble members, or the PSF shapes and sizes used for the TENOR-SAXS analysis. The distribution function was also changed (presented in Fig.\,\ref{fig:p_distribution_braid}). The \textit{baseline parameters} were Diamond Light Source B21 beamline characteristics \cite{Cowieson2020_B21} with the scattering ensemble featuring 11 normally distributed members, with  $\left<R_g\right>=5\,\text{nm}$. The observables extracted at different conditions are shown in Fig.\,\ref{fig:sensitiVT} - \ref{fig:sensitiVT2}, using the parameters detailed in Tab.\ \ref{tab:beam_simul_param}. The horizontal span of each braid of dashed curves indicates the uncertainty of the deduction of the ensemble's polydispersity from the given an observation and the scattering. In other words, for a given observable value the uncertainty in $V$ depends on the width of the braid for the variability under examination. 

We are here

As expected, we find that the radius of gyration, $R_0$, hardly plays a role in the spherical and Guinier form-factors (Fig. \ref{fig:sensitiVT}). Figure \ref{fig:sensitiVT3} shows that the instrumental PSF plays a negligivble role, confirming the main claim of the TENOR-SAXS method.The instrumental PSFs checked cover both synchrotron and lab-based typical PSFs. In Fig.~\ref{fig:sensitiVT1} we show that the number of different radii (distribution resolution) does not affect the method's performance.
In contrast to the ensemble and instrumental properties, there is a choice of variable used in the application of the method, the smearing PSF quartet size (as detailed in Sec.~\ref{sec:extraction2D}. Figure \ref{fig:sensitiVT2} shows that this has a larger effect than the other parameters discussed in this section. However, we should bear in mind that this is a controlled parameter, and as long as it is kept the same in the application on the actual measurements and the ground truth calibration, the dependence does not play a role. We did find that in noisy simulations, the smearing PSF is better when it is as large as possible (but within the Guinier region).


\textit{Signal strengh evaluation}: To obtain a representative estimate of realistic signal-to-noise levels, and to interpret the needed number of collected photons, we consider the Diamond Light Source bioSAXS beamline B21 as an example for a challenging biological case. The beamline delivers  $\sim (2-4)\times 10^{12}$~photons per second at the sample \cite{Cowieson2020_B21,Diamond_B21_Website}. 
In absolute units, the macroscopic scattered-photon yield per unit solid angle is  $\dot{N}_\Omega \;\approx\; \Phi \, D \, I(q)$, where $\Phi$ is the incident flux (photons $\mathrm{s}^{-1}$) at the sample, $D\approx 1\,\mathrm{mm}$ is the illuminated path length, and $I(q)$ is the absolute scattering intensity expressed in $\mathrm{cm}^{-1}$. 
Typical values of the peak scattering intensity of proteins are $I(0)/c\sim(1-50)\times 10^{-2} ~\mathrm{cm}^{-2}\mathrm{mg}^{-1}$ \cite{mylonas2007}. Typical protein concentrations are 0.1-10 $\mathrm{mg}\,\mathrm{cm}^{-3}$ \cite{mohammed2024}. Let us take $I(0)\approx1\,\mathrm{cm}^{-1}$, which obtains a scattering yield rate per solid angle of
\[
\dot{N}_{\Omega} \;\approx\;  3\times 10^{11} \ \text{photons\,sec}^{-1}\text{sr}^{-1}.
\]
Converting to $q^{2}$-space (using the small-angle approximation
 $q \approx k\theta$) employs the constant Jacobian $d\Omega/d(q^{2})
    = \pi k^{-2}
    = \lambda^{2}/(4\pi)$. At x-ray photons of 13~keV ($\lambda = 0.095\,\mathrm{nm}$), this factor evaluates to $7.18\times 10^{-4}\,\mathrm{sr}\,\mathrm{nm}^{2}$.  
Therefore, our estimated forward scattering rate converts to
\[
\dot N_{q^{2}} 
    \approx 2\times10^{8}\,
    \mathrm{s^{-1}\,nm^{2}}.
\]
Practically, this means that at such cases, one may attain the $10^{9} \text{ photons\,nm}^{2} $ scattered photon rate needed for efficient extraction of the ensemble variance using the TENOR-SAXS approach, by collecting less than ten frames (seconds).
\begin{table}
    \centering
    \footnotesize
    \begin{tabular}{l|l|p{5cm}}\hline
         Parameter& Baseline value& Inspected span\\\hline\hline
         Wavelength&  0.1 nm& \\\hline
         Sample-detector distance&  360 cm& \\\hline
         Detector pixel size&  $70\mu m\times70\mu m$& \\\hline
         Beam divergence  (slit size in pixels)&  $3\times15$& \\ \hline\hline
 Mean radius of gyration& $R_0=5nm$&2.5-7nm\\\hline
Ensemble polydispersity& 11&7-23\\\hline
 Size distribution function& Normal&Boltzmann, Log-normal, Schulz, Triangular, Uniform\\\hline
% & &\\
% & &\\
 Form-factor& $\phi''\in[-0.02,0.06]$&also spherical shell, solid sphere, thin rod, Gaussian chain\\\hline
 Smearing PSFs (pixels of $4\sigma$ kernel)& [(11,17)\,(3,5)]&[(11,17)\,(3,5)]- [(101,117)\,(113,115)]\\\hline
 Relative ensemble $R_g$ variance ($V$)& &0-0.3\\
    \end{tabular}
    \caption{Parameters used in the simulations. The top part lists the resolution based on Diamond Light Source B21 beamline's configuration \cite{Cowieson2020_B21}. The bottom part details the ensemble parameters and the method's variables.}
    \label{tab:beam_simul_param}
\end{table}

\begin{figure}
\includegraphics[width=0.9\textwidth]{observ_rg2.eps}
\caption{Sensitivity to different ensemble and method parameters in the extraction of the observables.\label{fig:sensitiVT}}
\end{figure}
\begin{figure}
\includegraphics[width=0.9\textwidth]{observ_PSF0.eps}
\caption{Sensitivity to different instrumental PSF size (expressed in its footprint's number of detector pixels) in the extraction of the observables.\label{fig:sensitiVT3}}
\end{figure}
\begin{figure}
\includegraphics[width=0.9\textwidth]{observ_populat2.eps}
\caption{Sensitivity to different ensemble and method parameters in the extraction of the observables.\label{fig:sensitiVT1}}
\end{figure}
\begin{figure}
\includegraphics[width=0.9\textwidth]{observ_PX2.eps}
\caption{Sensitivity to different ensemble and method parameters in the extraction of the observables.\label{fig:sensitiVT2}}
\end{figure}
\printbibliography
\end{document}
